\documentclass{article}
\usepackage[margin=0.6in]{geometry}
\usepackage{bm}
\usepackage{amsmath,amsthm,amssymb,amsfonts, fancyhdr, color, comment, environ}
\usepackage{xcolor}
\usepackage{mdframed}
\usepackage[shortlabels]{enumitem}
\usepackage{indentfirst}
\usepackage{hyperref}
\usepackage{graphicx}
\usepackage{tikzsymbols}
\usepackage{tcolorbox}
\usepackage{mathpazo}
\usepackage[object=vectorian]{pgfornament}
\hypersetup{
    colorlinks=true,
    linkcolor=blue,
    filecolor=magenta,      
    urlcolor=blue,
}

\renewcommand{\a}{\texttt{a}}
\renewcommand{\b}{\texttt{b}}
\renewcommand{\c}{\texttt{c}}
\renewcommand{\d}{\texttt{d}}

\newcommand{\0}{\texttt{0}}
\newcommand{\1}{\texttt{1}}
\newcommand{\2}{\texttt{2}}
\newcommand{\h}{\texttt{\#}}

\usepackage{tikz}
\usepackage{float}
\usetikzlibrary{automata, positioning, arrows}
\tikzset{
%->, % makes the edges directed
>=stealth, % makes the arrow heads bold
node distance=3cm, % specifies the minimum distance between two nodes. Change if necessary.
every picture/.style=thick,
every state/.style={thick, fill=gray!10}, % sets the properties for each ’state’ node
initial text=$ $, % sets the text that appears on the start arrow
}

\newcommand{\sectionline}{%
  \noindent
  \begin{center}
  {\color{black}
    \resizebox{0.5\linewidth}{1ex}
    {{%
    {\begin{tikzpicture}
    \node  (C) at (0,0) {};
    \node (D) at (9,0) {};
    \path (C) to [ornament=85] (D);
    \end{tikzpicture}}}}}%
    \end{center}
  }


\pagestyle{fancy}


\newenvironment{problem2}[2][Problem]
    { \begin{mdframed}[backgroundcolor=white] \textbf{#1 #2} \\}
    {  \end{mdframed}}

\newenvironment{problem}[5][Problem]
    { \begin{tcolorbox}[colback=white, title= {#1 #2 (#3): #4 (#5 points)}]}
    {  \end{tcolorbox}}

% Define solution environment
\newenvironment{solution}
    {\textit{Proof:}}
    {}

\renewcommand{\qed}{\quad\qedsymbol}

% prevent line break in inline mode
\binoppenalty=\maxdimen
\relpenalty=\maxdimen

\newcommand{\examinfo}[6]{
\rhead{\Huge{\boxed{\textsf{\textbf{#1}}}}}
%\rhead{\inclraphics[scale=.2]{bracuLogo.jpeg}}
\lhead{#2\\ #3} 
\chead{\textsc{#4 \\Total marks: #5}}
}

% default heading
\def\coursecode{Spring 2024}
\def\coursetitle{Automata and Computability}
\def\examtype{Midterm}
\def\marks{40}
\def\duration{80}

\newcommand{\setexamheading}{
\examinfo{\problemset}{\coursecode}{\coursetitle}{\examtype}{\marks}{\duration}
}

\newenvironment{examinstructions}{
    \begin{tcolorbox}
    \begin{center}
}{
    \end{center}
    \end{tcolorbox}
}

%%%%%%%%%%%%%%%%%%%%%%%%%%%%%%%%%%%%%%%%%%%%%%%%%%%%%%%%%%%%%%%%%%%%%%%%%%%%%%%%%%%%


\def\coursetitle{Automata and Computability}
\def\examtype{Assignment}
\def\marks{25}
\def\problemset{CSE331}

%%%%%%%%%%%%%%%%%%%%%%%%%%%%%%%%%%%%%%%%%%%%%

\setexamheading

\begin{document}


\begin{examinstructions}
There are a total of \(n\) problems. You have to solve all of them.
\end{examinstructions}



\begin{problem}{1}{CO5}{Nonregular Language}{25}
Use the pumping lemma to \textbf{demonstrate} that \(L_1\) and \(L_2\) are not regular.

\begin{enumerate}[(a)]
\item \( L_1 = \{ w \in \{\0, \1\}^* : w = \0^{n!} \text{ where } n \geq 0 \} \) (5 points)
\item \( L_2 = \{ w \in \{\0, \1\}^* : w = \0^a\1^b\1^c \0^d\text{ where } a+b=c+d \text{ and }  a,b,c,d \geq 0 \  \} \) (5 points)
\end{enumerate}
\end{problem}




\sectionline
\begin{enumerate}[(a)]
\item Assume for the sake of contradiction that \(L_1\) is regular. Then let \(p\) be the pumping length for \(L_1\). Now we take the string 
\[ w = \0^{p!} \in L_1. \]
Then the length of \(w\) is \(|w| = p! \geq p\). So \(w\) can be split into \(xyz\) such that \(|y| > 0\), \(|xy| \leq p\), and \(xy^iz \in L_1\) for each \(i \geq 0\). \(y\) consists of only \(\0\)s, so 
\[ xy^i z = \0^{p! + (i-1) |y| } \]
Then, for \(i=2\), \(xy^2z\) will be 
\[ xy^2z = xyyz = \0^{p! + |y|} .  \]
Now, \(|y| \le p < p \cdot p! \), hence, 
\[ p! < p! + |y| < p!  + p \cdot p! = p! \left(1+p\right) = (p+1)! \]
So \(p! < p! + |y| < (p+1)!\), and the length of \(xy^2z\) is strictly between two consecutive factorials. Hence, it cannot be a factorial. 
Thus we get a contradiction! Hence, \(L_1\) is not a regular language. 

% \item[(a)] \textbf{(Alternate Solution)} Assume for the sake of contradiction that \(L_1\) is regular. Then its complement language \(\overline L_1\) is also regular. 
% \[ \overline L_1 = \{ w \in \{\0, \1\}^* : w = \0^{m} \text{ where } m \ne n! \text{ for any } n \ge 0 \} . \]
% Let \(p\) be the pumping length of \(\overline L_1\). Now we take the string 
% \[w = \0 ^{p \cdot p!} \in \overline L_1 . \]
% Then the length of \(w\) is \(|w| = p \cdot p! \geq p\). So \(w\) can be split into \(xyz\) such that \(|y| > 0\), \(|xy| \leq p\), and \(xy^iz \in L_1\) for each \(i \geq 0\). \(y\) consists of only \(\0\)s, so 
% \[ xy^i z = \0^{p \cdot p! + (i-1) |y| } . \]
% Now, since \(0 < |y| \le p\), we have that \(|y|\) divides \( p! \). Now, choosing \(i = 1 + \frac{p!}{|y|}\), we have 
% \[ xy^i z = \0^{p \cdot p! + (i-1) |y| } 
% = \0^{p \cdot p! + \frac{p!}{|y|} |y| } 
% = \0 ^{p \cdot p! + p! }
% = \0 ^{(p+1)! } \notin \overline L_1 . \]
% Therefore, \(\overline L_1\) is not a regular language. Thus we get a contradiction, and hence \(L_1\) cannot be regular as well. 

\item Assume for the sake of contradiction that \(L_2\) is regular. Then let \(p\) be the pumping length for \(L_2\). Now we take the string 
\[ w = \0^{p} \1 \1 \0^p \in L_2. \]
Then the length of \(w\) is \(|w| = 2p+2 \geq p\). So \(w\) can be split into \(xyz\) such that \(|y| > 0\), \(|xy| \leq p\), and \(xy^iz \in L_2\) for each \(i \geq 0\). \(|xy| \le p\), so \(y\) consists of only \(\0\)s, so 
\[ xy^i z = \0^{p + (i-1) |y| } \1 \1 \0^p . \]
We choose \(i=4\), so that 
\[ xy^4 z = \0^{p + 3 |y| } \1 \1 \0^p . \]
Since this is in \(L_2\), one can write it as \(\0^a\1^b\1^c \0^d\) for \(a+b=c+d\). By equating 
\[ \0^{p + 3 |y| } \1 \1 \0^p = \0^a\1^b\1^c \0^d ,  \]
we get \(a= p + 3 |y|\), \(d = p\) and \(b+c=2\). So \(c-b \le 2\). Furthermore, 
\[ c-b = a-d =  3 |y| \ge 3 ,  \]
as \(|y| \ge 1\). So we get \(c-b \le 2\) and \(c-b \ge 3\), which is a contradiction! Therefore, \(L_2\) is not regular. 

\item [(b)] \textbf{(Alternate Solutiion)} Assume for the sake of contradiction that \(L_2\) is regular. Then let \(p\) be the pumping length for \(L_2\). Now we take the string 
\[ w = \0^{p} \1 ^p \1 ^p \0^p \in L_2. \]
Then the length of \(w\) is \(|w| = 4p \geq p\). So \(w\) can be split into \(xyz\) such that \(|y| > 0\), \(|xy| \leq p\), and \(xy^iz \in L_2\) for each \(i \geq 0\). \(|xy| \le p\), so \(y\) consists of only \(\0\)s, so 
\[ xy^i z = \0^{p + (i-1) |y| } \1 ^p \1 ^p \0^p . \]
We choose \(i=2p+2\), so that 
\[ xy^{2p+2} z = \0^{p + (2p+1) |y| } \1 ^p \1 ^p \0^p . \]
Since this is in \(L_2\), one can write it as \(\0^a\1^b\1^c \0^d\) for \(a+b=c+d\). By equating 
\[ \0^{p + (2p+1) |y| } \1 ^p \1 ^p \0^p = \0^a\1^b\1^c \0^d ,  \]
we get \(a= p + (2p+1) |y|\), \(d = p\) and \(b+c=2p\). So \(c-b \le 2p\). Furthermore, 
\[ c-b = a-d =  (2p+1) |y| \ge 2p+1 ,  \]
as \(|y| \ge 1\). So we get \(c-b \le 2p\) and \(c-b \ge 2p+1\), which is a contradiction! Therefore, \(L_2\) is not regular. 

\end{enumerate}

\end{document}